\documentclass[11pt,a4paper]{article}
\usepackage[T1]{fontenc}
\usepackage[utf8]{inputenc}
\usepackage[english]{babel}
\usepackage{amsmath,amssymb}
\usepackage[margin=2cm]{geometry}
\setlength{\parindent}{0pt}
\setlength{\emergencystretch}{3em}
\allowdisplaybreaks

\title{Studies of functions: where the derivatives are positive, negative or zero (37)}
\author{}
\date{}
\begin{document}
\maketitle
Every statement below is true. Each says where the first or second derivative of a function is positive, negative or zero: the intervals of monotonicity or of concavity, and the stationary points, printed in an exercise collection or derived from its printed extrema. A comparison counts as false where one of its sides has no value, and a derivative has a value only where the function is defined on both sides of the point. The engine must prove them all.

\begin{enumerate}
\item $\displaystyle \forall x \in \mathbb{R}\ \left(\frac{d}{dx}\left( x \cdot (2 - 3x)^{3} \right) > 0 \iff x < \frac{1}{6}\right)$
\item $\displaystyle \forall x \in \mathbb{R}\ \left(\frac{d}{dx}\left( x \cdot (2 - 3x)^{3} \right) < 0 \iff \frac{1}{6} < x < \frac{2}{3} \lor x > \frac{2}{3}\right)$
\item $\displaystyle \forall x \in \mathbb{R}\ \left(\frac{d}{dx}\left( \frac{x^{2} + 3}{x - 5} \right) > 0 \iff x < 5 - 2\sqrt{7} \lor x > 5 + 2\sqrt{7}\right)$
\item $\displaystyle \forall x \in \mathbb{R}\ \left(\frac{d}{dx}\left( \frac{x^{2} + 3}{x - 5} \right) < 0 \iff 5 - 2\sqrt{7} < x < 5 \lor 5 < x < 5 + 2\sqrt{7}\right)$
\item $\displaystyle \forall x \in \mathbb{R}\ \left(\frac{d}{dx}\left( \frac{x^{2}}{x^{2} + 3x - 3} \right) < 0 \iff 0 < x < \frac{-3 + \sqrt{21}}{2} \lor \frac{-3 + \sqrt{21}}{2} < x < 2\right)$
\item $\displaystyle \forall x \in \mathbb{R}\ \left(\frac{d}{dx}\left( \frac{x^{2} - 3}{x + 2} \right) > 0 \iff x < -3 \lor x > -1\right)$
\item $\displaystyle \forall x \in \mathbb{R}\ \left(\frac{d}{dx}\left( \frac{x^{2} - 3}{x + 2} \right) < 0 \iff -3 < x < -2 \lor -2 < x < -1\right)$
\item $\displaystyle \forall x \in \mathbb{R}\ \left(\frac{d}{dx}\left( x\ln^{3} x \right) > 0 \iff e^{-3} < x < 1 \lor x > 1\right)$
\item $\displaystyle \forall x \in \mathbb{R}\ \left(\frac{d}{dx}\left( x\ln^{3} x \right) < 0 \iff 0 < x < e^{-3}\right)$
\item $\displaystyle \forall x \in \mathbb{R}\ \left(\frac{d}{dx}\left( x^{4} + 2x^{3} - 2x^{2} \right) > 0 \iff -2 < x < 0 \lor x > \frac{1}{2}\right)$
\item $\displaystyle \forall x \in \mathbb{R}\ \left(\frac{d}{dx}\left( \frac{x^{5}}{5} - \frac{x^{4}}{2} - x^{3} \right) < 0 \iff -1 < x < 0 \lor 0 < x < 3\right)$
\item $\displaystyle \forall x \in \mathbb{R}\ \left(\frac{d}{dx}\left( \frac{125}{729}(x - 2)^{3}(x + 1)^{2} \right) > 0 \iff x < -1 \lor \frac{1}{5} < x < 2 \lor x > 2\right)$
\item $\displaystyle \forall x \in \mathbb{R}\ \left(\frac{d}{dx}\left( \frac{x^{2} - x}{1 + x^{2}} \right) > 0 \iff x < -\sqrt{2} - 1 \lor x > \sqrt{2} - 1\right)$
\item $\displaystyle \forall x \in \mathbb{R}\ \left(\frac{d}{dx}\left( x^{2} + 2 + \frac{x^{3}}{4 - x} \right) > 0 \iff 0 < x < 4 \lor 4 < x < 8\right)$
\item $\displaystyle \forall x \in \mathbb{R}\ \left(\frac{d}{dx}\left( \frac{1}{x\sqrt{1 + x}} \right) > 0 \iff -1 < x < -\frac{2}{3}\right)$
\item $\displaystyle \forall x \in \mathbb{R}\ \left(\frac{d}{dx}\left( \frac{1}{x\sqrt{1 + x}} \right) < 0 \iff -\frac{2}{3} < x < 0 \lor x > 0\right)$
\item $\displaystyle \forall x \in \mathbb{R}\ \left(0 \le x \le 2\pi \implies \left(\frac{d}{dx}\left( \sin x(1 + \cos x) \right) > 0 \iff 0 \le x < \frac{\pi}{3} \lor \frac{5\pi}{3} < x \le 2\pi\right)\right)$
\item $\displaystyle \forall x \in \mathbb{R}\ \left(\frac{d}{dx}\left( \arctan(x^{4} - x^{3}) \right) < 0 \iff x < 0 \lor 0 < x < \frac{3}{4}\right)$
\item $\displaystyle \forall x \in \mathbb{R}\ \left(\frac{d}{dx}\left( -x^{3} + 6x^{2} \right) > 0 \iff 0 < x < 4\right)$
\item $\displaystyle \forall x \in \mathbb{R}\ \left(\frac{d^{2}}{dx^{2}}\left( 2x^{3} - 3x^{2} - 12x + 4 \right) > 0 \iff x > \frac{1}{2}\right)$
\item $\displaystyle \forall x \in \mathbb{R}\ \left(\frac{d^{2}}{dx^{2}}\left( x^{4} - 8x^{2} - 3 \right) > 0 \iff x < -\frac{2}{\sqrt{3}} \lor x > \frac{2}{\sqrt{3}}\right)$
\item $\displaystyle \forall x \in \mathbb{R}\ \left(\frac{d^{2}}{dx^{2}}\left( \ln(x^{2} + 1) \right) > 0 \iff -1 < x < 1\right)$
\item $\displaystyle \forall x \in \mathbb{R}\ \left(\frac{d^{2}}{dx^{2}}\left( x \cdot e^{-x} \right) > 0 \iff x > 2\right)$
\item $\displaystyle \forall x \in \mathbb{R}\ \left(\frac{d^{2}}{dx^{2}}\left( x^{3}e^{x} \right) > 0 \iff -3 - \sqrt{3} < x < -3 + \sqrt{3} \lor x > 0\right)$
\item $\displaystyle \forall x \in \mathbb{R}\ \left(\frac{d^{2}}{dx^{2}}\left( x^{3}e^{x} \right) < 0 \iff x < -3 - \sqrt{3} \lor -3 + \sqrt{3} < x < 0\right)$
\item $\displaystyle \forall x \in \mathbb{R}\ \left(\frac{d^{2}}{dx^{2}}\left( x^{5}e^{-x} \right) < 0 \iff x < 0 \lor 5 - \sqrt{5} < x < 5 + \sqrt{5}\right)$
\item $\displaystyle \forall x \in \mathbb{R}\ \left(\frac{d}{dx}\left( \sin(3x) \right) = 0 \iff \exists k \in \mathbb{Z}\ \left(x = \frac{\pi}{6} + \frac{k\pi}{3}\right)\right)$
\item $\displaystyle \forall x \in \mathbb{R}\ \left(\frac{d}{dx}\left( x^{3}e^{x} \right) > 0 \iff -3 < x < 0 \lor x > 0\right)$
\item $\displaystyle \forall x \in \mathbb{R}\ \left(\frac{d}{dx}\left( \ln(x^{2} - 1) \right) > 0 \iff x > 1\right)$
\item $\displaystyle \forall x \in \mathbb{R}\ \lnot\left(\frac{d}{dx}\left( \ln x \right) < 0\right)$
\item $\displaystyle \forall x \in \mathbb{R}\ \left(\frac{d}{dx}\left( \sqrt{x} \right) > 0 \iff x > 0\right)$
\item $\displaystyle \forall x \in \mathbb{R}\ \left(\frac{d}{dx}\left( \arcsin x \right) > 1 \iff -1 < x < 0 \lor 0 < x < 1\right)$
\item $\displaystyle \forall x \in \mathbb{R}\ \left(\frac{d^{2}}{dx^{2}}\left( \sqrt{x} \right) < 0 \iff x > 0\right)$
\item $\displaystyle \forall x \in \mathbb{R}\ \left(\frac{d}{dx}\left( \arctan x \right) > 0\right)$
\item $\displaystyle \forall x \in \mathbb{R}\ \left(\frac{d}{dx}\left( \frac{d}{dx}\left( x^{3} \right) \right) = 6x\right)$
\item $\displaystyle \forall x \in \mathbb{R} \setminus \{0\}\ \left(\frac{d}{dx}\left( \ln|x| \right) = \frac{1}{x}\right)$
\item $\displaystyle \forall x \in \mathbb{R} \setminus \{0\}\ \left(\frac{d}{dx}\left( x|x| \right) = 2|x|\right)$
\end{enumerate}
\end{document}
